\section{Routing, rooted models, and minimizing path systems}
\label{sec:routing}

The minor constructions in the next two sections require two kinds of
control over their connecting paths. We must insert rooted clique models
inside districts while preserving the paths' prescribed ends. We must also
bound the list-coloring cost of deleting long path systems. We first
establish the packing and replacement tools, then prove the sparsity bounds
for minimum-length systems and show how to combine two families sharing
a target set.

All paths in this section are simple; a path with one vertex is allowed when
its two prescribed ends coincide.  The length of a path is its number of
edges.  A family of paths is vertex-disjoint unless explicitly permitted to
share prescribed starting vertices.  When a path is oriented from its source
to its target, its first and last vertices refer to that orientation.

\subsection{Packing paths and rooting a clique minor}

We begin with the vertex-capacity form of Menger's theorem
\cite{menger1927}, proved here by the integral max-flow/min-cut argument
of Ford and Fulkerson \cite{fordfulkerson1956}.

\begin{lemma}[Vertex path packing]\label{lem:packing}
Let $A,B$ be disjoint vertex sets in a graph $P$, and let $j\ge1$ be an
integer.  Suppose that for every $X\subseteq V(P)$ with $|X|<j$, the graph
$P-X$ contains a path from $A\setminus X$ to $B\setminus X$.
Then there are $j$ pairwise vertex-disjoint paths from distinct vertices of
$A$ to distinct vertices of $B$.  Each path can be chosen to meet $A\cup B$
only at its two ends.
\end{lemma}

\begin{proof}
Replace each vertex $v$ by two vertices $v^-,v^+$ and an arc
$v^-v^+$ of capacity one.  For each edge $uv$, add arcs $u^+v^-$ and
$v^+u^-$, both of capacity $j$.  Add a source $s$ with arcs $sa^-$ of
capacity $j$ for $a\in A$, and a sink $z$ with arcs $b^+z$ of capacity $j$
for $b\in B$.

Start with zero flow.  Whenever the residual directed graph has an
$s$--$z$ path, augment by one unit along it.  Residual capacities remain
integral.  Stop upon reaching value $j$, or if no residual $s$--$z$ path
exists.  In the latter case let $R$ be the vertices reachable from $s$ in
the residual graph.  Every original arc from $R$ to its complement is
saturated, and every original arc entering $R$ from its complement has
zero flow.  Flow conservation therefore shows that the capacity of the
outgoing cut equals the current flow value.

If this value is less than $j$, no capacity-$j$ arc crosses the cut.
Its outgoing arcs are consequently unit arcs $v^-v^+$, for a set $X$ of
fewer than $j$ original vertices.  Any $A$--$B$ path in $P-X$ would give
an $s$--$z$ path avoiding those cut arcs, which is impossible.  This
contradicts the hypothesis.  Thus an integral flow of value $j$ is reached.

Decompose that flow into unit $s$--$z$ paths and directed cycles, and
discard the cycles.  Each original vertex is used by at most one of the
unit paths because its unit arc has capacity one.  The corresponding
paths in $P$ are therefore vertex-disjoint.  On each such path, take its
first visit to $B$ and its last preceding visit to $A$, and retain the
segment between them.  This gives the asserted endpoint convention
without changing the number or disjointness of the paths.
\end{proof}

\begin{lemma}[Two paths from every source]\label{lem:double-packing}
Let $A,B$ be disjoint vertex sets in a graph $P$, where $r=|A|\ge1$.
If $P$ is $2r$-connected and $|B|\ge2r$, there are two paths starting at
each vertex of $A$, ending at $2r$ distinct vertices of $B$, with no
intersections other than the two paths sharing their prescribed start.
Their internal vertices lie outside $A\cup B$.
\end{lemma}

\begin{proof}
Replace each vertex of $A$ by two copies.  For every original edge, join
all copies of one endpoint to all copies of the other, treating a vertex
outside $A$ as having a single copy.  Write $A^*$ for the $2r$ source
copies.  The target set $B$ is unchanged.

Let $X^*$ be a set of fewer than $2r$ vertices of this new graph.
At least one source copy and at least one target survive.  In the
original graph delete the set $X$ of vertices all of whose copies were
deleted.  Then $|X|\le |X^*|<2r$, so the surviving original source and
target are connected in $P-X$.  A simple path between them lifts to a
path in the copy graph minus $X^*$ by choosing a surviving copy at each
vertex.  Thus $X^*$ cannot separate $A^*$ from $B$.

Lemma~\ref{lem:packing} now gives $2r$ vertex-disjoint paths from $A^*$
to $B$.  Every source copy is an endpoint of one path, so no source copy
is internal to another path.  Identifying the two copies of each source
therefore introduces exactly the permitted common starts and no other
overlaps.  The endpoint convention in Lemma~\ref{lem:packing} also
excludes internal vertices in $B$.
\end{proof}

A stronger rooted-minor theorem of Kawarabayashi is recorded in
\cite[Lemma~5.14]{DelcourtPostle2021}. We include a proof of the form
needed here.

\begin{lemma}[Rooting a large clique minor]\label{lem:rooting}
Let $k\ge1$.  If $P$ is $2k$-connected and $h(P)\ge4k$, then any given
$k$ distinct vertices are roots of a $K_k$ model: its bags can be indexed
by the given vertices so that each bag contains its indexing root.
\end{lemma}

\begin{proof}
We use two successive path packings. The first sends two paths from each
root towards the given clique model and leaves some model bags untouched.
Its path--bag unions may overlap, but each vertex lies in at most two of
them. This bound will certify a second packing to untouched bags, producing
the disjoint rooted model.

Let $R$ be the prescribed root set.  Start with $4k$ disjoint bags of a
clique model.  At most $k$ of these bags meet $R$, so choose $3k$ bags
disjoint from $R$ and call them eligible.  Choose one representative in
each eligible bag.  Lemma~\ref{lem:double-packing}, applied to these
representatives as targets, gives two paths from each root, with their
$2k$ ends in distinct eligible bags.

Among all such double systems, now allowing arbitrary endpoints within
the eligible bags but still requiring distinct end bags, choose one of
minimum total length.  No path meets an eligible bag that is not an end
bag of the system: its first such visit would allow that path to be
truncated and assigned the unused bag, strictly reducing total length.
Consequently at least $k$ eligible bags are untouched by every path.
Choose $k$ of them and contract each to a single target vertex; let $T$
be the resulting set of $k$ targets.

For each of the $2k$ paths of the double system, take its union with its
own end bag.  These $2k$ connected sets are unchanged by the contractions
just made.  Every vertex belongs to at most two of them.  Indeed, a
vertex outside $R$ belongs to at most one path and at most one eligible
bag, whereas a root belongs to its two paths and to no eligible bag.
Paths may pass through other used end bags; this observation already
allows for those intersections.

In the contracted graph, delete fewer than $k$ vertices.  They meet
fewer than $2k$ of the connected sets just described, so at least one
such set survives intact, including its root.  At least one vertex of
$T$ also survives.  The surviving set has an edge to every surviving
vertex of $T$, because its end bag was adjacent to every one of the
contracted eligible bags.  Hence these deletions cannot separate $R$
from $T$.

Apply Lemma~\ref{lem:packing} with $j=k$.  All $k$ roots and all $k$
targets are used as endpoints, and none is internal on another path.
Expand each contracted target bag and add to it the vertices of its
root-to-target path.  The resulting bags are connected and disjoint.
They remain pairwise adjacent through the original edges between the
target bags, and each contains its assigned root.
\end{proof}

\subsection{Wovenness and replacement inside districts}

Wovenness records the simultaneous rooted-model and path property used in
Postle's framework \cite[Section~3.3]{Postle2020}; see also
\cite[Definition~5.12]{DelcourtPostle2021}. We specify the endpoint
coincidences because the recursive construction uses them.

\begin{definition}\label{def:woven}
For integers $a\ge1$ and $b\ge0$, a graph $J$ is \emph{$(a,b)$-woven}
if the following holds.  Prescribe $a$ distinct roots
$R=\{r_1,\ldots,r_a\}$ and at most $b$ pairs
$(x_1,y_1),\ldots,(x_p,y_p)$ of vertices of $J$.  The endpoint sets
$\{x_i,y_i\}$ must be disjoint for different $i$, but $x_i=y_i$ is
allowed, and endpoints may belong to $R$.  Then there are a $K_a$ model
with bags $M_1,\ldots,M_a$ satisfying $r_j\in M_j$, and pairwise
vertex-disjoint $x_i$--$y_i$ paths $Q_i$, such that
\[
 V(Q_i)\cap\bigcup_{j=1}^a M_j
 \ \subseteq\ R\cap\{x_i,y_i\}\qquad(1\le i\le p).
\]
For a coincident pair the path is the single prescribed vertex.
Because the model bags are disjoint and contain all the roots, any
allowed intersection at $r_j$ belongs to the designated bag $M_j$.
\end{definition}

We use the woven-subgraph replacement argument of
\cite[Lemma~5.15]{DelcourtPostle2021}, recording its intersection
conclusions at the first and last district visits of each path.

\begin{lemma}[Rerouting through a woven district]\label{lem:rerouting}
Let $J$ be an $(a,b)$-woven subgraph of a graph $P$, and prescribe $a$
distinct roots in $J$.  Let $Q_1,\ldots,Q_q$ be vertex-disjoint paths in $P$,
where $q\le b$.  Orient each path arbitrarily.  For every path meeting
$J$, let $x_i,y_i$ be its first and last vertices in $J$.
There are vertex-disjoint paths $Q'_1,\ldots,Q'_q$ with the same
respective endpoints, and a model rooted at the prescribed roots in $J$,
such that the following properties hold:
\begin{enumerate}[label=(\roman*)]
\item every vertex added to a path belongs to $J$; more precisely,
      $V(Q'_i)\setminus V(J)\subseteq V(Q_i)\setminus V(J)$;
\item a path originally disjoint from $J$ is unchanged;
\item the intersection of $Q'_i$ with the model is contained in the
      prescribed roots among $x_i,y_i$.
\end{enumerate}
The construction can be applied successively in disjoint districts.
If some prescribed roots are endpoints of the original path system,
the bags for these roots meet the resulting paths only at their own
roots, including after all subsequent replacements in other districts.
\end{lemma}

\begin{proof}
The pairs $(x_i,y_i)$ are admissible in Definition~\ref{def:woven}:
their endpoint sets are disjoint across paths, and a path meeting $J$
only once gives a permitted coincident pair.  Apply wovenness to all
these pairs.  On each path meeting $J$, replace the entire segment from
$x_i$ to $y_i$ by the path supplied inside $J$.  Its retained prefix and
suffix have no vertices in $J$ except their attachment endpoints.
They are disjoint from the interiors of all the replacement paths.
The retained pieces of different original paths are themselves
disjoint, and the replacement paths are disjoint by wovenness.
The resulting concatenations are therefore simple, vertex-disjoint
paths with the required endpoints and intersections.  Notice that a
replaced segment may have made excursions outside $J$; those excursions
are removed, explaining the inclusion rather than equality in (i).

For successive replacements, the new vertices at a later step lie in
a district disjoint from every earlier district and its model.  Thus
no new intersection with an earlier model is introduced.  A prescribed
root that is an original path endpoint remains that endpoint at every
step.  In its own district it is a first or last visit of that path,
and no other path contains it.  Its bag therefore has precisely the
allowed path intersection.
\end{proof}

\begin{remark}\label{routing:padding-roots}
The last conclusion does not require every prescribed root to be a
path endpoint.  An auxiliary root can be a first or last visit of an
unrelated path, and its bag can then meet that path at this root.
If only bags whose roots are designated path endpoints are retained,
discarding the auxiliary bags removes this issue.  A root that was
strictly between the first and last district visits need not remain
on that path after rerouting.
\end{remark}

Fix the positive integer $\gamma$ from Theorem~\ref{thm:grouping}, and set
\begin{equation}\label{routing:K0}
 K_0=10^5\gamma.
\end{equation}

\begin{lemma}[Small districts are woven]\label{lem:small-woven}
For all sufficiently large integers $t$, every $(K_0t)$-connected graph
$J$ of order at most $t^{1.06}$ is $(s,3t)$-woven whenever
\[
 1\le s\le 2\left\lceil\frac{t}{(\log t)^2}\right\rceil.
\]
The cutoff is absolute and uniform in $s$ and $J$.  No minor-exclusion
hypothesis on $J$ is required.
\end{lemma}

\begin{proof}
We set aside the prescribed roots and pair endpoints, then split the
remaining vertices into one part containing an unrooted clique model
and a second, still highly connected part in which to join its bags to
the roots and route all prescribed pairs.

Fix an admissible root set and a system of at most $3t$ pairs.
Let $Z$ be the union of the roots and all pair endpoints.  For sufficiently
large $t$ the displayed upper bound on $s$ is at most $t$, so $|Z|\le7t$.
Put $k'=\lfloor K_0t/2\rfloor$.  Deleting $Z$ leaves a $k'$-connected
graph: deleting a further set of fewer than $k'$ vertices still deletes
fewer than $K_0t$ vertices in total, and the order condition also follows
from $|V(J)|>K_0t$.  Every vertex of $J$, including each vertex of $Z$,
has at least $K_0t-|Z|\ge k'$ neighbors outside $Z$.

Independently place each vertex of $J-Z$ in $A$ with probability $1/1000$,
and write $B'=V(J)\setminus(Z\cup A)$.  We show that with probability
tending to one both of the following hold:
\begin{align}
 |N_J(v)\cap A|&\ge k'/2000,
 &|N_J(v)\cap B'|&\ge k'/2 &&(v\in V(J)),
 \label{routing:split-degrees}\\
 J[B']&\text{ is }\lfloor k'/32\rfloor\text{-connected}.
 \label{routing:split-connectivity}
\end{align}
For each vertex of $J$, fix $k'$ of its neighbors outside $Z$.
The number of these neighbors placed in $A$ is a sum of $k'$
independent indicators with mean $k'/1000$.  By
Lemma~\ref{lem:concentration}, the probability that it is below
$k'/2000$ or above $k'/2$ is $\exp(-\Omega(k'))$.
A union bound over at most $t^{1.06}$ vertices proves
\eqref{routing:split-degrees} with probability tending to one.

For \eqref{routing:split-connectivity}, form in advance the cut-test
matchings of Lemma~\ref{lem:cut-test} for $J-Z$, with connectivity
parameter $k'$ and order bound $n=t^{1.06}$.  The number of tests is at most
\[
 \exp\!\bigl(O(t^{0.06}(\log t)^2)\bigr).
\]
A test matching has at most $k'/4$ endpoints.  The expected number of
its endpoints in $A$ is at most $k'/4000$, whereas the cut test allows
fewer than $k'/32$.  This leaves a fixed positive multiple of $k'$ as
a deviation margin.  Lemma~\ref{lem:concentration}, applied to at most
$k'/4$ independent indicators, bounds each failure probability by
$\exp(-\Omega(k'))$.  Their union has probability tending to zero.
On the degree event \eqref{routing:split-degrees}, every vertex of
$J-Z$ has at least $k'/2$ neighbors outside $A$.  Thus all hypotheses
of Lemma~\ref{lem:cut-test} hold simultaneously with positive probability,
and it gives \eqref{routing:split-connectivity}.  Fix such a split.

The graph $J[A]$ is nonempty and has minimum degree at least $k'/2000$;
hence its density is at least $k'/4000=\Omega(t)$.  Uniformly for the
allowed values of $s\ge2$,
\[
 s\log(2s)=O\!\left(\frac{t}{\log t}\right)=o(t).
\]
Lemma~\ref{lem:coarse-density} therefore implies that $J[A]$ contains
a $K_s$ model, for all sufficiently large $t$.  If $s=1$, a single
vertex of $A$ supplies the model.  Denote the model bags by
$M_1^0,\ldots,M_s^0$ and select a port $p_j\in M_j^0$ from each bag.

Choose distinct representatives in $B'$ for all the following roles:
each port, each root, and each of the two endpoints of every
noncoincident prescribed pair.  A representative is required to be a
neighbor of the vertex it represents.  When a root is also a pair
endpoint, its two roles receive different representatives.  There are
at most $2s+6t\le8t$ roles.  By \eqref{routing:split-degrees}, each has
at least $k'/2$ available neighbors in $B'$, so greedy distinct selection
is possible.

Partition the representatives into groups of size two: the port and
root representatives for each $j$, and the two endpoint representatives
for each noncoincident pair.  By \eqref{routing:K0}, for sufficiently
large $t$,
\[
 \left\lfloor\frac{k'}{32}\right\rfloor\ge \gamma(2s+6t).
\]
Theorem~\ref{thm:grouping} in $J[B']$ gives disjoint connected sets for
these groups.  Write $C_j$ for the set joining the representatives of
$p_j$ and $r_j$.  The set
\[
 M_j=M_j^0\cup C_j\cup\{r_j\}
\]
is connected, using the two representative attachment edges.  These
sets are disjoint and remain pairwise adjacent through the original
model on $A$.

In each other grouping set choose a path joining its two representatives,
and extend it by their attachment edges to the original pair endpoints.
Use the trivial path for every coincident pair.  The nontrivial paths
are disjoint inside $B'$, all their original endpoints lie in $Z$, and
different pairs have disjoint endpoint sets.  The model bags contain
no vertices of $Z$ except their own roots.  Consequently the only
model--path intersections are roots that are endpoints of the respective
pairs, exactly as in Definition~\ref{def:woven}.
\end{proof}

\subsection{Minimum-length systems with free target assignments}

The bounds below use the unpaired geodesic-path argument of Postle
\cite[Lemma~4.14]{Postle2020}, extended here to two paths at each source.
The freedom to reassign targets is essential: the proof shortens a
system by exchanging its path suffixes.

The graph in which a path system is minimized matters: all shortcut and
degree statements below concern edges of that same graph.  We shall use
two kinds of systems with disjoint source and target sets $A,B$:
\begin{enumerate}[label=\textup{(\arabic*)}]
\item two paths from each vertex of $A$, ending at distinct vertices of
      $B$, with intersections allowed only at their prescribed repeated
      starts;
\item a specified number $q$ of vertex-disjoint paths from arbitrary
      distinct vertices of $A$ to arbitrary distinct vertices of $B$.
\end{enumerate}
In either case there is no prescribed pairing between sources and targets.
For type~\textup{(1)}, the total number of paths is $q=2|A|$.

\begin{lemma}[Bounds from unpaired path minimization]
\label{lem:path-minimization}
Let $A,B$ be disjoint vertex sets in a graph $P$.  Suppose that a system
of one of the two types above is feasible, and choose one of minimum
total length.  Let $q$ be its number of paths and $U$ its vertex union.
Then:
\begin{enumerate}[label=(\roman*)]
\item the paths have no internal vertices in $A\cup B$;
\item each vertex of $V(P)\setminus U$ has at most three neighbors
      on any one path;
\item $P[U\setminus A]$ is $2q$-degenerate for type~\textup{(1)},
      and $P[U]$ is $2q$-degenerate for type~\textup{(2)}.
\end{enumerate}
\end{lemma}

\begin{proof}
An internal visit to $B$ would allow truncation at that visit, strictly
shortening a path.  Its new target cannot be the target of another path,
by disjointness.  For type~\textup{(2)}, an internal visit to $A$ likewise
allows deletion of the initial segment, making that visit the new source.
For type~\textup{(1)}, all sources are already prescribed starts, and
the allowed-intersection rule excludes any internal visit to $A$.
This proves (i).

For (ii), let $v\notin U$ have at least four neighbors on a path
$x_0x_1\cdots x_l$, ordered from source to target.  Its first and last
neighbor indices $i<j$ satisfy $j-i\ge3$.  If $v\notin A\cup B$,
replace the segment from $x_i$ to $x_j$ by $x_i v x_j$.  This saves at
least one edge and introduces no intersection with another path.
If $v\in B$, truncate the path after $x_i$ and use $v$ as its target;
here $i+1\le l-2<l$.  If $v\in A$, the system must have type~\textup{(2)},
since all type~\textup{(1)} sources belong to $U$.  Start instead at $v$
and follow the suffix from $x_j$ to $x_l$; its length is
$1+l-j\le l-2$.  These are valid systems with new free endpoint choices,
contradicting minimality in every case.

To prove (iii), put $U_0=U\setminus A$ for type~\textup{(1)} and
$U_0=U$ for type~\textup{(2)}.  Suppose that some nonempty induced
subgraph $P[S]$, $S\subseteq U_0$, has minimum degree greater than $2q$.
For every path meeting $S$, mark its first two vertices in $S$, or its
only such vertex if there is just one.  At most $2q$ vertices are marked.
Each vertex of $S$ belongs to exactly one path: the only possible
overlaps were the repeated starts, which were removed in type~\textup{(1)}.

For a path $Q_i$ meeting $S$, let $u_i$ be its first vertex in $S$.
Its degree in $P[S]$ is greater than $2q$, so it has a neighbor $v_i$
in $S$ which is not marked.  This vertex belongs to some path $Q_j$
and is at least its third vertex in $S$.  Direct an arrow from $i$ to
$j$, recording the edge $u_i v_i$.  Every participating path has one
outgoing arrow, so these arrows contain a directed cycle, possibly a loop.

For each arrow $i\to j$ on this cycle, form a new path from the prefix
of $Q_i$ ending at $u_i$, the edge $u_i v_i$, and the suffix of $Q_j$
starting at $v_i$.  Leave all paths outside the cycle unchanged.
On each cycle path $Q_j$, the retained prefix ends at $u_j$, while the
retained suffix begins at the incoming arrow's vertex, at least the
third vertex of $S$ on $Q_j$.  The two retained pieces are disjoint and
separated by at least two original path edges.  Pieces lying on
different original paths are disjoint except for the permitted repeated
starts.  Thus the new concatenations are simple and satisfy precisely
the original disjointness conditions.

Every original source prefix is retained, so the required sources and
their multiplicities are unchanged.  This remains true if both paths
with the same source belong to the directed cycle: only their suffixes
are exchanged, and their shared initial vertex remains their only
intersection.  The target suffixes are permuted, preserving distinct
targets.  Such a permutation is allowed because no pairing was prescribed.

If the cycle has $c$ paths, at least $2c$ original edges are removed
between their retained prefixes and suffixes, and precisely $c$ joining
edges are inserted.  The total length decreases by at least $c>0$.
This final contradiction proves (iii).
\end{proof}

\subsection{Combining two doubled families}

\begin{lemma}[Combining paths by pairing temporary sources]
\label{lem:combine}
Let $Z,D,V(H)$ be pairwise disjoint vertex sets in a graph $P$, and let
$m\ge1$.  Suppose there are:
\begin{enumerate}[label=(\roman*)]
\item two paths from each vertex of $Z$ to distinct vertices of $D$,
      with all target ends distinct, with intersections only at their
      prescribed repeated starts, and avoiding $V(H)$;
\item $2m$ vertex-disjoint paths from distinct vertices of $H$ to
      distinct vertices of $D$, avoiding $Z$ and meeting $H$ only at
      their starting vertices.
\end{enumerate}
Then, using only vertices of these two families, one can obtain
$|Z|+m$ vertex-disjoint paths to distinct vertices of $D$, one starting
at each vertex of $Z$ and one starting at each of $m$ distinct vertices
of $H$.  These paths meet $D$ and $H$ only at their appropriate ends.
Moreover, after any partition of the $2m$ second-family starts into $m$
pairs, the $m$ selected starts in $H$ can be required to include exactly
one vertex from each pair.  No unselected second-family start is visited
by any of the resulting paths.
\end{lemma}

\begin{proof}
First truncate every input path at its first visit to $D$ if necessary.
Within either family this preserves all distinctness and disjointness
requirements.  Work in the graph consisting only of the vertices and
edges of the two resulting path families; ambient edges are unnecessary.
For each prescribed pair of second-family starts, identify its two
vertices to form one temporary source.  The paths of that family have
no other vertices in $H$, and the first family avoids $H$, so each
original path still gives a path after identification.

Write $N=|Z|+m$.  There are now $N$ sources, namely $Z$ together with
the $m$ temporary sources, and there are two given routes from each
source to $D$.  Across these $2N$ routes, every vertex has congestion
at most two, where congestion counts the number of routes containing
the vertex.  At a vertex of $Z$ there are only its two first-family
routes, since the second family avoids $Z$.  At a temporary source
there are only its two second-family routes, since the first family
avoids $H$.  Elsewhere there is at most one route from each family.
This includes vertices of $D$, whose ends are distinct within each
family but need not be distinct between the two families.

Deleting fewer than $N$ vertices meets fewer than $2N$ routes.
Some entire source-to-$D$ route therefore survives.  By
Lemma~\ref{lem:packing}, there are $N$ disjoint paths from the full
source set to $D$, with no internal source or target vertices.
Every source is used, since the source set has cardinality $N$.

Each temporary source occurs on exactly one packed path and only as
its starting endpoint.  Lift the first edge of that path to whichever
one of its two original starts supplied this edge, and keep the rest
of the path unchanged.  This selects exactly one original start from
the pair.  No switching between the two original starts is needed,
because the identified vertex was not internal.  Different temporary
sources came from disjoint pairs, so all the lifted paths remain
vertex-disjoint.

The only vertices of $H$ in the graph of the original families were
the $2m$ second-family starts.  All their images were source vertices
and could not be internal in the packing.  Thus the lifted paths meet
$H$ exactly at the $m$ selected starts; all unselected starts are
absent.  The target-end convention survives lifting, and every used
vertex belongs to an original path family, as required.
\end{proof}
